\documentclass[12pt]{article}
\usepackage[utf8]{inputenc}
\usepackage{ctex} % 支持中文排版
\usepackage{geometry}
\geometry{a4paper, margin=1in}
\usepackage{amsmath}
\usepackage{amssymb}
\usepackage{setspace}

\begin{document}

\begin{center}
    {\Large \textbf{Joey Practice 1}} \\
\end{center}

\vspace{0.5cm}

\textbf{1. 从数据中找规律 (Finding Patterns in Data)}

观察下面的表格，其中 $x$ 代表输入值，$y$ 代表输出值。\\
(Observe the table below, where $x$ represents the input and $y$ represents the output.)

\begin{table}[ht]
\centering
\begin{tabular}{|c|c|}
\hline
$x$ (输入值/Input) & $y$ (输出值/Output) \\
\hline
1 & 8 \\
2 & 11 \\
3 & 14 \\
4 & 17 \\
5 & 20 \\
\hline
\end{tabular}
\end{table}

a) 当 $x$ 每次增加 1 时，$y$ 增加了多少？（这通常叫做变化率或斜率）\\
(As $x$ increases by 1 each time, how much does $y$ increase by? This is often called the rate of change or slope.)
\vspace{1.5cm}

b) 试着往前推算一下，如果 $x = 0$ 时，$y$ 应该是多少？（这通常叫做初始值或截距）\\
(Try working backwards: what would the value of $y$ be when $x = 0$? This is often called the initial value or y-intercept.)
\vspace{1.5cm}

c) 根据上面找出的规律，请写出能够描述这个规律的线性方程（格式通常为 $y = mx + b$）。\\
(Based on the pattern you found, please write the linear equation that describes it, usually in the form $y = mx + b$.)
\vspace{2cm}

d) 利用你刚才写出的方程算一下，当 $x = 45$ 时，$y$ 是多少？（请写出计算过程）\\
(Using the equation you just wrote, calculate the value of $y$ when $x = 45$. Please show your work.)
\vspace{3cm}

\newpage

\textbf{2. 坐标系里的直线 (Lines in the Coordinate System)}

已知一个线性方程：(Given the linear equation:) $y = -2x + 6$

a) 请自己选择几个连续的 $x$ 取值（包括正数和 0），算出对应的 $y$ 值，填在下面的表格里。\\
(Choose a few consecutive integer values for $x$ (including positive numbers and 0), calculate the corresponding $y$ values, and complete the table below.)

\vspace{0.3cm}
\begin{table}[ht]
\centering
\begin{tabular}{|c|c|c|}
\hline
$x$ & $y$ & 坐标点 (Coordinate Point) $(x, y)$ \\
\hline
0 & 6 & $(0, 6)$ \\
\hline
\raisebox{0pt}[15pt][10pt]{}  & & \\
\hline
\raisebox{0pt}[15pt][10pt]{}  & & \\
\hline
\raisebox{0pt}[15pt][10pt]{}  & & \\
\hline
\end{tabular}
\end{table}
\vspace{0.3cm}

b) 在坐标纸上（或自己画一个坐标系），标出表格里的四个点，并用尺子把它们连成一条直线（记得在直线两边画上箭头表示无限延伸）。\\
(On a grid or your own coordinate plane, plot the four points from your table and use a ruler to connect them into a straight line. Remember to draw arrows on both ends to show it extends infinitely.)
\vspace{6cm}

c) 观察你画的这条直线，它和 $y$ 轴交在哪一点？它是从左往右上升的还是下降的？这和方程里的哪个数字有关系？\\
(Look at the line you drew. At what point does it cross the $y$-axis? Is it rising or falling from left to right? Which number in the equation is this related to?)
\vspace{3cm}

\newpage

\textbf{3. 动手解方程与不等式 (Solving Equations and Inequalities)}

a) 请写出详细的步骤，求出下面方程里 $x$ 的值：\\
(Please show detailed steps to solve for $x$ in the equation below:)
\[ 5x - 8 = 2x + 13 \]
\vspace{5cm}

b) 请写出求解下面不等式的每一个步骤：\\
(Please show every step to solve the inequality below:)
\[ -4x + 7 \ge 23 \]
\vspace{4cm}

c) 画一条数轴，并在上面表示出上面这道不等式的解。（注意要不要涂黑圆圈，以及箭头的方向）\\
(Draw a number line and graph the solution to the inequality above. Pay attention to whether the circle should be closed or open, and the direction of the arrow.)
\vspace{4cm}

\newpage

\textbf{4. 生活中的数学：手机套餐 (Math in Real Life: Cell Phone Plans)}

某通信公司提供一款手机上网套餐，计费规则是：\\
每个月的“固定基础费”是 \$15.00。除此之外，每用 1 GB 的流量，要多付 \$2.50。\\
假设 $x$ 是这个月用的流量（单位：GB），$C$ 是这个月的总账单费用（单位：美元）。\\
(A telecommunications company offers a mobile internet plan with the following billing rules: \\
The monthly "fixed base fee" is \$15.00. In addition, you pay an extra \$2.50 for every 1 GB of data used. \\
Let $x$ be the data used this month (in GB) and $C$ be the total monthly bill (in dollars).)

a) 请用代数来帮忙，写出总费用 $C$ 的计算方程式。\\
(Please use algebra to help write the equation to calculate the total cost $C$.)
\vspace{2.5cm}

b) 如果小明这个月用了 8 GB 流量，他的总账单会是多少钱？\\
(If Xiaoming used 8 GB of data this month, how much would his total bill be?)
\vspace{3cm}

c) 小华这个月收到了一份 \$47.50 的账单。请列出一个方程，算算小华这个月用了多少 GB 流量，并写出求解步骤。\\
(Xiaohua received a bill for \$47.50 this month. Please set up an equation to figure out how many GB of data Xiaohua used this month, and show the steps to solve it.)
\vspace{6cm}

\newpage

\textbf{5. 附加挑战题 (Bonus Challenges)}

\textbf{挑战A：双重边界 (Challenge A: Double Boundaries)}

解一个“被夹在中间”的不等式，求出 $x$ 的取值范围，并写出解题步骤。\\
(Solve the following compound inequality to find the range of values for $x$, and show your steps.)
\[ -12 < 3x - 6 \le 15 \]
\vspace{4cm}

然后，试着用 \textbf{区间表示法 (Interval Notation)} 把它写出来（比如考虑用圆括号 () 还是方括号 [] 来表示大大小小）。\\
(Then, try writing it using \textbf{Interval Notation}. Think about whether to use parentheses () or brackets [] to represent the boundaries.)
\vspace{2.5cm}


\textbf{挑战B：两条直线的交点 (Challenge B: Intersection of Two Lines)}

这里有两个方程：(Here are two equations:)\\
方程 A (Equation A)：$y = x + 3$ \\
方程 B (Equation B)：$x + y = 7$

a) 方程 B 现在看起来不太像我们熟悉的 $y = ...$ 这种形式。请把它变一下形，变成以 $y$ 等于什么的样子。\\
(Equation B doesn't look like the familiar $y = ...$ form right now. Please rearrange it so it's solved for $y$.)
\vspace{2cm}

b) 如果把这两条直线画在同一个坐标系里，它们分别长什么样？（你可以在演草纸上画一下）\\
(If you graph these two lines on the same coordinate plane, what would they look like? You can sketch them on scratch paper.)
\vspace{4.5cm}

c) 它们相交在哪个点？你能找出这交点的精确坐标吗？\\
(At what point do they intersect? Can you find the exact coordinates of this intersection?)
\vspace{1.5cm}

d) 算出来的这个 $x$ 和 $y$ 的值，如果放回原来的方程 A 和方程 B 里，会怎么样？这有什么特别的意义吗？\\
(If you plug the calculated $x$ and $y$ values back into original Equations A and B, what happens? What special significance does this have?)
\vspace{3cm}

\end{document}